% Claim statement file for row claim_3_15 -- "AddingInterventionNodes".
% Auto-generated stub by scaffold/solve_chapter.py at row start. The
% manager agent (or a worker it dispatches) fills the body below with the
% claim's statement(s) from the lecture notes -- statement only, no proof
% (proofs live in the sibling `claim_3_15_proof_*.tex` file).
%
% This file is a sibling of `main.tex` inside `<subsection>/tex/`. The
% `subfiles` package lets it render both standalone and as part of
% `main.tex`. Note: `main.tex` does NOT \subfile this statement file -- the
% sibling `_proof_` file restates the statement above the proof, so
% including both in `main.tex` would be redundant. This file exists for
% standalone reading / grep convenience.

\documentclass[main]{subfiles}

\begin{document}
% ref: claim_3_15 (statement)
% title: AddingInterventionNodes
\def\rowref{claim\_3\_15}
\phantomsection\label{claim_3_15}

\begin{Lem}[Adding intervention nodes commutes with disjoint node-splitting hard interventions]\label{adding-intervention-nodes-commutes-with-disjoint-node-splitting-hard-interventions}
    Let $G = (J, V, E, L)$ be a CADMG (def \ref{def-cdmg-types}, item~i., equivalently a CDMG of def \ref{def-cdmg} that is acyclic in the sense of def \ref{def-acylic}); throughout we use the notation of def \ref{not-cdmg} and recall the typing and axioms of def \ref{def-cdmg}, namely $J \cap V = \emptyset$, $E \subseteq (J \cup V) \times V$, $L \subseteq V \times V$, $L$ irreflexive ($(v_1, v_2) \in L \Rightarrow v_1 \neq v_2$), and $L$ symmetric ($(v_1, v_2) \in L \Leftrightarrow (v_2, v_1) \in L$). We rely throughout on the node-splitting hard-intervention (SWIG) construction $\swig$ of def \ref{def:G_node-splitting_intervention} and the intervention-node extension $\doit(I_{\cdot})$ of def \ref{def:cdmg_intervention_nodes}. Let
    \[ W_1 \subseteq V \quad\text{and}\quad W_2 \subseteq J \cup V \]
    be two subsets of nodes of $G$ (both automatically finite, since $J \cup V$ is finite by def \ref{def-cdmg}; this matches the project's Lean encoding of $W_1$ and $W_2$ as \texttt{Finset}s). The typings $W_1 \subseteq V$ (strictly tighter than the ambient $J \cup V$, since $W_1$ must lie on the output side) and $W_2 \subseteq J \cup V$ (overlap with $J$ is permitted) are dictated by the preconditions of the two constructors involved: $\swig$ acts only on output nodes (def \ref{def:G_node-splitting_intervention}, ``Let $W \subseteq V$''), while $\doit(I_{\cdot})$ admits an arbitrary subset of the carrier (def \ref{def:cdmg_intervention_nodes}, ``Let $W \subseteq J \cup V$''). The two subsets are subject to the disjointness side condition
    \[ W_1 \cap W_2 \;=\; \emptyset. \]
    All three quantifiers ``for every CADMG $G$'', ``for every $W_1 \subseteq V$'', and ``for every $W_2 \subseteq J \cup V$ with $W_1 \cap W_2 = \emptyset$'' are read \emph{universally}; in particular the (vacuously disjoint) corner cases $W_1 = W_2 = \emptyset$, $W_1 = \emptyset$ alone, and $W_2 = \emptyset$ alone are admitted, and so are the following two corner cases governed by the $I_j := j$ convention of def \ref{def:cdmg_intervention_nodes}:
    \begin{itemize}
        \item $W_2 \subseteq J$ (so $W_2 \sm J = \emptyset$): the extension $G_{\doit(I_{W_2})} = G$ is the identity by def \ref{def:cdmg_intervention_nodes} (no fresh intervention nodes and no new edges are added, since the $I_j := j$ shorthand for $j \in J \cap W_2 = W_2$ is purely notational); the displayed equality below then degenerates to $G_{\swig(W_1)} = G_{\swig(W_1)}$, which is trivially true. The $W_2 \subseteq J$ case is admitted by the statement (read literally per def \ref{def:cdmg_intervention_nodes}, not naively as ``$|W_2|$ fresh nodes added'').
        \item $W_2 \cap J \neq \emptyset$ and $W_2 \cap V \neq \emptyset$ (the partial case): the extension $G_{\doit(I_{W_2})}$ adds fresh intervention nodes $\lC I_w \st w \in W_2 \sm J \rC$ and edges $\lC (I_w, w) \st w \in W_2 \sm J \rC$ only for the $V$-elements $W_2 \sm J = W_2 \cap V$ of $W_2$, and acts as the identity on the $J$-elements $W_2 \cap J$ (the $I_j := j$ shorthand contributes no new node and no new edge there). Sub-case admitted; the equality below holds with the $W_2 \sm J$ parametrisation visible in the directed-edge component (per def \ref{def:cdmg_intervention_nodes}, item~iii.).
    \end{itemize}

    \emph{Reading of ``CADMG'' in the LN's prose.}
    The LN's literal text opens with ``Let $G = (J, V, E, L)$ be a CADMG'' and refers in the prose to ``the CADMG that arises from $\ldots$'' for both sides of the displayed equality. Both sides are produced by applying $\swig$ (def \ref{def:G_node-splitting_intervention}) and $\doit(I_{\cdot})$ (def \ref{def:cdmg_intervention_nodes}) in some order to $G$. The constructor of def \ref{def:G_node-splitting_intervention} requires the CADMG hypothesis on its input -- on $G$ for the LHS's inner $\swig(W_1)$, and on $G_{\doit(I_{W_2})}$ for the RHS's outer $\swig(W_1)$. Acyclicity of $G_{\doit(I_{W_2})}$ under the CADMG hypothesis on $G$ is established by \ref{claim_3_13} (\emph{AcyclicHardInterventionTopologicalOrder} in the SWIG-analog presentation of this section, more precisely the remark right after def \ref{def:cdmg_intervention_nodes} in the LN, which states that acyclicity of $G$ implies acyclicity of $G_{\doit(I_W)}$ for every $W \subseteq J \cup V$). The constructor of def \ref{def:cdmg_intervention_nodes} only requires its input to be a CDMG (not a CADMG); since both $G$ and $G_{\swig(W_1)}$ are CDMGs (the latter by the CDMG-axiom check at the end of def \ref{def:G_node-splitting_intervention}), the inner $\doit(I_{W_2})$ on either side is defined regardless of acyclicity. Hence both sides of the displayed equality are CADMGs. At the underlying CDMG-type level, the asserted equality below is first an equality of CDMGs (def \ref{def-cdmg}) -- componentwise on the four components $(J, V, E, L)$ -- and the upgrade to an equality of CADMGs is automatic once both sides are known to be acyclic CDMGs (acyclicity is a propositional condition on a fixed CDMG, transported along the CDMG-equality).

    \emph{Disambiguation of the $\swig(\cdot)$ and $\doit(I_{\cdot})$ notation.}
    The only $\swig$ in this lemma is $\swig(W_1)$ with $W_1 \subseteq V$, and the only $\doit$ in this lemma is $\doit(I_{W_2})$ with the $I$-tagged argument $I_{W_2}$ (not the bare-node-set argument $W_2$). In particular, the hard-intervention reading $G_{\doit(W)}$ of def \ref{def:G_hard_intervention} (which would consume a bare node set $W$ and produce a CDMG with $J_{\doit(W)} := J \cup W$ etc.) is \emph{not} invoked here; nor is any SWIG-of-$I$-tagged-set notation. Concretely:
    \begin{itemize}
        \item $G_{\swig(W_1)}$ denotes the \emph{node-splitting hard intervention} on $W_1$ (def \ref{def:G_node-splitting_intervention}, which requires $G$ to be a CADMG -- matching the standing hypothesis on $G$, and giving the SWIG construction directly via four set-builders rather than via an intermediate node-split). Its components are
              $J_{\swig(W_1)} := J \dcup W_1^i$,
              $V_{\swig(W_1)} := (V \sm W_1) \dcup W_1^o$,
              $E_{\swig(W_1)} := \lC (v_1^i, v_2^o) \st (v_1, v_2) \in E \rC$, and
              $L_{\swig(W_1)} := \lC (v_1^o, v_2^o) \st (v_1, v_2) \in L \rC$,
              where $W_1^o := \lC w^o \st w \in W_1 \rC$ and $W_1^i := \lC w^i \st w \in W_1 \rC$ are the two type-level tagged copies of $W_1$ introduced in def \ref{def:G_node-splitting_intervention} (so $W_1^o \cap W_1^i = \emptyset$ and both $W_1^o$ and $W_1^i$ are disjoint from every element of $J \cup V$, the latter automatically by the tagged construction), and the shorthand $v^o := v^i := v$ is used for $v \in J \cup (V \sm W_1)$ inside the set-builders for $E_{\swig(W_1)}$ and $L_{\swig(W_1)}$ (def \ref{def:G_node-splitting_intervention}, items~i.--iv.). The input-side tagged copy $W_1^i$ sits in $J_{\swig(W_1)}$ and the output-side tagged copy $W_1^o$ sits in $V_{\swig(W_1)}$.
        \item $G_{\doit(I_{W_2})}$ denotes the \emph{extended CDMG} obtained by adding intervention nodes for $W_2$ (def \ref{def:cdmg_intervention_nodes}). Its components are
              $J_{\doit(I_{W_2})} := J \cup \lC I_w \st w \in W_2 \sm J \rC$,
              $V_{\doit(I_{W_2})} := V$,
              $E_{\doit(I_{W_2})} := E \cup \lC (I_w, w) \st w \in W_2 \sm J \rC$, and
              $L_{\doit(I_{W_2})} := L$,
              where $\lC I_w \st w \in W_2 \sm J \rC$ is the family of fresh intervention symbols introduced by def \ref{def:cdmg_intervention_nodes} (type-disjoint from $J \cup V$ and pairwise type-disjoint), and the notational shorthand $I_j := j$ for $j \in J \cap W_2$ means no new node is introduced and no new edge is added on the $J \cap W_2$ branch. The output-node set $V$ and the bidirected-edge set $L$ are untouched.
    \end{itemize}

    \emph{Well-typedness of the iterated operations on both sides.}
    Each iterated operation appearing in the displayed equality below is well-typed per its constructor's precondition:
    \begin{itemize}
        \item \emph{LHS, $\lp G_{\swig(W_1)} \rp_{\doit(I_{W_2})}$.} The inner $G_{\swig(W_1)}$ is defined per def \ref{def:G_node-splitting_intervention} from the standing $W_1 \subseteq V$ and the CADMG hypothesis on $G$. By items~i.\ and ii.\ of def \ref{def:G_node-splitting_intervention}, the carrier of $G_{\swig(W_1)}$ is
              \[ J_{\swig(W_1)} \cup V_{\swig(W_1)} \;=\; \lp J \dcup W_1^i \rp \;\cup\; \lp (V \sm W_1) \dcup W_1^o \rp, \]
              with the canonical inclusion (the ``unsplit'' map) $\iota_{\swig} \colon J \cup (V \sm W_1) \hookrightarrow J_{\swig(W_1)} \cup V_{\swig(W_1)}$ acting as $v \mapsto v$ -- and the elements of $W_1$ themselves are \emph{not} in the image of $\iota_{\swig}$, having been split into the tagged copies $W_1^o, W_1^i$. The disjointness $W_1 \cap W_2 = \emptyset$ together with $W_2 \subseteq J \cup V$ gives $W_2 \subseteq J \cup (V \sm W_1)$, so $\iota_{\swig}$ restricts to an inclusion $W_2 \hookrightarrow J_{\swig(W_1)} \cup V_{\swig(W_1)}$ identifying $W_2$ with its image $\iota_{\swig}[W_2]$ as the unsplit-untagged copy of $W_2$ inside the SWIG carrier. Hence the outer $\doit(I_{W_2})$ is defined per def \ref{def:cdmg_intervention_nodes} applied to $G_{\swig(W_1)}$ with the subset $\iota_{\swig}[W_2] \subseteq J_{\swig(W_1)} \cup V_{\swig(W_1)}$ (def \ref{def:cdmg_intervention_nodes} only requires its input graph to be a CDMG, and $G_{\swig(W_1)}$ is a CDMG by the axiom-check at the end of def \ref{def:G_node-splitting_intervention}). Under this identification, $\iota_{\swig}[W_2] \cap J_{\swig(W_1)} = W_2 \cap J$ (the $J$-elements of $W_2$ remain in $J$, with no node addition under the $I_j := j$ convention) and $\iota_{\swig}[W_2] \sm J_{\swig(W_1)} = W_2 \sm J \subseteq V \sm W_1$ (the $V$-elements of $W_2$ sit on the untagged $V \sm W_1$ piece, with one fresh $I_w$ added for each).
        \item \emph{RHS, $\lp G_{\doit(I_{W_2})} \rp_{\swig(W_1)}$.} The inner $G_{\doit(I_{W_2})}$ is defined per def \ref{def:cdmg_intervention_nodes} from the standing $W_2 \subseteq J \cup V$ on the CDMG underlying the CADMG $G$. By items~i.\ and ii.\ of def \ref{def:cdmg_intervention_nodes}, the carrier of $G_{\doit(I_{W_2})}$ is
              \[ J_{\doit(I_{W_2})} \cup V_{\doit(I_{W_2})} \;=\; \lp J \cup \lC I_w \st w \in W_2 \sm J \rC \rp \;\cup\; V, \]
              with the canonical inclusion (the ``unsplit'' map) $\iota_{\doit} \colon J \cup V \hookrightarrow J_{\doit(I_{W_2})} \cup V_{\doit(I_{W_2})}$ acting as $v \mapsto v$ for all $v \in J \cup V$ (every original node of $G$ embeds untagged into the extended carrier; only the fresh symbols $I_w$ for $w \in W_2 \sm J$ are new). Since $W_1 \subseteq V = V_{\doit(I_{W_2})}$, the inclusion $\iota_{\doit}$ restricts to an inclusion $W_1 \hookrightarrow V_{\doit(I_{W_2})}$ identifying $W_1$ with its image $\iota_{\doit}[W_1]$ on the output side. Hence the outer $\swig(W_1)$ is defined per def \ref{def:G_node-splitting_intervention} applied to $G_{\doit(I_{W_2})}$ with the subset $\iota_{\doit}[W_1] \subseteq V_{\doit(I_{W_2})}$. The CADMG-hypothesis precondition of def \ref{def:G_node-splitting_intervention} on $G_{\doit(I_{W_2})}$ is met because $G_{\doit(I_{W_2})}$ is acyclic whenever $G$ is, by the LN's remark immediately following def \ref{def:cdmg_intervention_nodes} (already noted under ``Reading of `CADMG' '' above).
    \end{itemize}

    \emph{Distinct carriers and the canonical relabelling identifying them.}
    The LHS and RHS of the displayed equality live in formally distinct iterated-tagged-sum carriers:
    \begin{itemize}
        \item the LHS carrier is $J_{\lp G_{\swig(W_1)} \rp_{\doit(I_{W_2})}} \cup V_{\lp G_{\swig(W_1)} \rp_{\doit(I_{W_2})}}$, obtained by first splitting $W_1 \subseteq V$ off into the tagged copies $W_1^o, W_1^i$ (def \ref{def:G_node-splitting_intervention}) and then adding fresh intervention symbols $\lC I_w \st w \in W_2 \sm J \rC$ on top (def \ref{def:cdmg_intervention_nodes}); concretely
              \[ J \cup W_1^i \cup \lC I_w \st w \in W_2 \sm J \rC \;\cup\; (V \sm W_1) \cup W_1^o, \]
              with all four unions automatically disjoint by the tagged-copy disjointness of $W_1^o, W_1^i$ from $J \cup V$ and from each other, and by the freshness of the intervention symbols $\lC I_w \st w \in W_2 \sm J \rC$ from $J \cup V$ and from $W_1^o, W_1^i$ (the last clause follows from the tagged construction together with $W_1 \cap W_2 = \emptyset$, which ensures no $I_w$ symbol is accidentally identified with any $w' \in W_1$);
        \item the RHS carrier is $J_{\lp G_{\doit(I_{W_2})} \rp_{\swig(W_1)}} \cup V_{\lp G_{\doit(I_{W_2})} \rp_{\swig(W_1)}}$, obtained by first adding fresh intervention symbols $\lC I_w \st w \in W_2 \sm J \rC$ (def \ref{def:cdmg_intervention_nodes}) and then splitting $W_1 \subseteq V$ off into the tagged copies $W_1^o, W_1^i$ (def \ref{def:G_node-splitting_intervention}); concretely
              \[ J \cup \lC I_w \st w \in W_2 \sm J \rC \cup W_1^i \;\cup\; (V \sm W_1) \cup W_1^o, \]
              again with all four unions automatically disjoint.
    \end{itemize}
    The two carriers are equal as plain sets of symbols (both are the union $J \cup (V \sm W_1) \cup W_1^o \cup W_1^i \cup \lC I_w \st w \in W_2 \sm J \rC$, with the $\cup$ between any two of the five pieces automatically disjoint), and the canonical relabelling identifying them is the obvious one: the unsplit copy of every $v \in J \cup V$ maps to itself on either side; the tagged copies $w^o \in W_1^o$ and $w^i \in W_1^i$ on the LHS (introduced by ``first split'') are identified with the tagged copies $w^o \in W_1^o$ and $w^i \in W_1^i$ on the RHS (introduced by ``last split'') via the identity on $W_1$; and the fresh intervention symbols $I_w$ for $w \in W_2 \sm J$ on the LHS (introduced by ``last extend'') are identified with the fresh intervention symbols $I_w$ for $w \in W_2 \sm J$ on the RHS (introduced by ``first extend'') via the identity on $W_2 \sm J$. The asserted equality of CADMGs below is read componentwise on the four components $(J, V, E, L)$ of def \ref{def-cdmg} \emph{up to this canonical relabelling}; equivalently, when each carrier is identified with the common set just displayed, the four components of the LHS and the four components of the RHS coincide as subsets of (the carrier, the carrier $\times$ carrier with appropriate input--output typing, and the output-side carrier squared).

    \emph{The conclusion.} With the well-typedness and carrier identifications spelled out above, the following equality of CADMGs holds:
    \[ \lp G_{\swig(W_1)} \rp_{\doit(I_{W_2})} \;=\; \lp G_{\doit(I_{W_2})} \rp_{\swig(W_1)}. \]
    Componentwise -- which is the unique unpacking the equality admits, since CDMGs are tuples $(J, V, E, L)$ of def \ref{def-cdmg} -- the equality unfolds into the four field-by-field identifications (under the canonical relabelling of the previous paragraph):
    \begin{align*}
        J_{\lp G_{\swig(W_1)} \rp_{\doit(I_{W_2})}} \;&=\; J_{\lp G_{\doit(I_{W_2})} \rp_{\swig(W_1)}} \;=\; J \;\cup\; W_1^i \;\cup\; \lC I_w \st w \in W_2 \sm J \rC, \\
        V_{\lp G_{\swig(W_1)} \rp_{\doit(I_{W_2})}} \;&=\; V_{\lp G_{\doit(I_{W_2})} \rp_{\swig(W_1)}} \;=\; (V \sm W_1) \;\cup\; W_1^o, \\
        E_{\lp G_{\swig(W_1)} \rp_{\doit(I_{W_2})}} \;&=\; E_{\lp G_{\doit(I_{W_2})} \rp_{\swig(W_1)}} \;=\; \lC (v_1^i, v_2^o) \st (v_1, v_2) \in E \rC \;\cup\; \lC (I_w, w) \st w \in W_2 \sm J \rC, \\
        L_{\lp G_{\swig(W_1)} \rp_{\doit(I_{W_2})}} \;&=\; L_{\lp G_{\doit(I_{W_2})} \rp_{\swig(W_1)}} \;=\; \lC (v_1^o, v_2^o) \st (v_1, v_2) \in L \rC,
    \end{align*}
    where on the right-hand sides the tagged shorthands $v_1^i$, $v_1^o$, $v_2^o$ for $v_1, v_2 \in J \cup V$ are read per def \ref{def:G_node-splitting_intervention}, items~i.--iv., i.e.\ $v_k^o := v_k^i := v_k$ for $v_k \in J \cup (V \sm W_1)$ (untagged) and $v_k^o \in W_1^o$, $v_k^i \in W_1^i$ are the tagged copies when $v_k \in W_1$; and where in the second set-builder of the $E$-component the symbol $w$ for $w \in W_2 \sm J$ refers to its unsplit-untagged image $\iota_{\swig}[w] = w \in V \sm W_1$ (well-defined since $w \in W_2 \sm J \subseteq V \sm W_1$ by $W_1 \cap W_2 = \emptyset$ and $W_2 \sm J \subseteq V$). The conjunctive unpacking of the displayed equality into these four field-by-field equalities is deferred to the proof.
\end{Lem}

\end{document}
